\documentclass[11pt,a4paper]{article}

\usepackage[T1]{fontenc}
\usepackage[utf8]{inputenc}
\usepackage{amsmath,amssymb}
\usepackage[margin=2.5cm]{geometry}
\usepackage{microtype}
\usepackage[hidelinks]{hyperref}

\setlength{\parskip}{0.6em}
\setlength{\parindent}{0pt}
\pagestyle{empty}

\title{\vspace{-1.5cm}\bfseries Quark/gluon jet discrimination in dijet photoproduction at the EIC}

\author{%
  Siddharth Narayan Singh$^{1}$,\quad
  R.~Aggarwal$^{2}$,\quad
  M.~Kaur$^{3}$
}

\date{}

\begin{document}
\maketitle
\thispagestyle{empty}
\vspace{-1.2cm}

{\small
$^{1}$Brandeis University, Waltham, USA \\
$^{2}$USAR and USBAS, Guru Gobind Singh Indraprastha University,
East Delhi Campus, 110092, India \\
$^{3}$Department of Physics, Panjab University, Chandigarh 160014, India;
Department of Physics, Amity University, Punjab, Mohali 140306, India \\[0.3em]
\texttt{siddharthsingh@brandeis.edu}
}

\vspace{0.8em}
\hrule
\vspace{1.0em}

Gluon jets are wider than quark jets, and quark/gluon tagging exploits that. We compare three observables in dijet
photoproduction: the integrated jet shape $\Psi(r=0.3)$, which measures how a
jet's transverse energy is spread in radius; subjet
multiplicity $n_{\mathrm{subjets}}$, and the iterated soft-drop multiplicity $n_{\mathrm{SD}}$, which
counts splittings passing a momentum-fraction cut. All three are computed on
exactly the same jets, from PYTHIA~8 events at $\sqrt{s} = 64$,
$105$, $141$ and $300$~GeV, covering the three EIC design points and a HERA
reference.

Even though there is an increasing preference for machine learning techniques,
theoretically grounded methods remain crucial for improving the interpretability
of results and addressing biases in the simulated or experimental data used for
training. There are many other complex observables that can be used to tag the
jet based on the same physical content, expressed in different functional forms.
This paper presents the first study of soft-drop multiplicity in dijet
photoproduction events and compares it with the integrated jet shape and subjet
multiplicity.

\end{document}
